\section{Quantitative limits of local Abel data}
\label{sec:abel-jet-stability}

The construction in \eqref{round26:eq:factorization} has a stronger
approximation property than agreement of the Abel zero order. Its
normalized Abel derivatives can approximate those of Liouville to any
prescribed finite accuracy, while a real pole still prevents the critical
energy estimate. This distinguishes an exact arithmetic identity from
an approximation to it. None of the results below excludes a zeta zero.

Fix $1/2<\beta<1$ and an integer $P\ge2$. For primes $p>P$, put
\[
 q_p=\tfrac12p^{\beta-1},\qquad
 S(t)=\sum_{P<p\le t}q_p,\qquad
 I_p=\lfloor S(p)\rfloor-\lfloor S(p)-q_p\rfloor.
\]
Set $f_P(p)=-1$ for $p\le P$, set $f_P(p)=2I_p-1$ otherwise,
and extend completely multiplicatively. Write
\[
 \mathcal S_P=\{p>P:I_p=1\},\qquad
 G_P(s)=\prod_{p\in\mathcal S_P}\frac{1+p^{-s}}{1-p^{-s}}.
\]
In particular $f_P(n)=\lambda(n)$ for every integer $n\le P$.
The finite cutoff and $\beta$ specify one fixed function; they do not
change as the scale of an energy estimate tends to infinity.

\begin{lemma}[An explicit Euler-factor bound]\label{lem:jet-euler-bound}
For every real $\sigma>\beta$, the logarithm specified by its Euler
series is holomorphic on $\R s>\beta$ and satisfies
\begin{equation}\label{eq:jet-delta}
 \sup_{\R s\ge\sigma}|\log G_P(s)|\le\Delta(P,\sigma),\qquad
 \Delta(P,\sigma)=
 \frac{P^{\beta-\sigma}}
 {(\sigma-\beta)(1-P^{-2\sigma})}.
\end{equation}
Consequently $|G_P(s)-1|\le e^{\Delta(P,\sigma)}-1$ in that half-plane.
\end{lemma}
\begin{proof}
The centered cumulative discrepancy has a sign:
\[
 E(t)=\sum_{P<p\le t}(I_p-q_p)=\lfloor S(t)\rfloor-S(t)\in(-1,0].
\]
For $T>P$, partial summation gives
\[
 \sum_{P<p\le T}(I_p-q_p)p^{-\sigma}
 =E(T)T^{-\sigma}
   +\sigma\int_P^T E(t)t^{-\sigma-1}\dd t\le0.
\]
It follows, by monotone convergence of the nonnegative sums, that
\begin{align*}
 \sum_{p\in\mathcal S_P}p^{-\sigma}
 &\le\tfrac12\sum_{p>P}p^{-(1+\sigma-\beta)}\\
 &\le\tfrac12\sum_{n>P}n^{-(1+\sigma-\beta)}
 \le\frac{P^{\beta-\sigma}}{2(\sigma-\beta)}.
\end{align*}
The last inequality compares a decreasing sum with its integral from
$P$ to infinity. Thus the following series converges absolutely and
locally uniformly for $\R s>\beta$:
\[
 \log G_P(s)=2\sum_{p\in\mathcal S_P}
             \sum_{\substack{j\ge1\\j\text{ odd}}}\frac{p^{-js}}j.
\]
For $\R s\ge\sigma$, its absolute value is at most
\[
 \frac{2}{1-P^{-2\sigma}}
       \sum_{p\in\mathcal S_P}p^{-\sigma}\le\Delta(P,\sigma).
\]
Exponentiating proves the final assertion.
\end{proof}

The quotient called $H_f$ in the preceding construction includes a
finite Euler factor $c_P$. Here it is convenient to cancel that common
factor and normalize the double zero at one. Define
\begin{equation}\label{eq:normalized-abel-jet}
 \mathcal J_\lambda(s)=\frac{\zeta(2s)}{((s-1)\zeta(s))^2},\qquad
 \mathcal J_P(s)=\mathcal J_\lambda(s)G_P(s).
\end{equation}
For $f=f_P$, this is precisely $c_P(s)H_f(s)/(s-1)^2$; for
$f=\lambda$, the analogous identity holds with the same cutoff.
Both functions have their removable values assigned at one.
\footnote{A finite Abel jet of order $K$ consists of these derivatives
at one for $0\le k\le K$. It specifies local data.}

\begin{theorem}[Quantitative approximation of every finite Abel jet]
\label{thm:finite-abel-jets}
Let $0<r<1-\beta$ and put
\[
 C_r=\frac{\zeta(2-2r)}{(1-2r)^2}.
\]
The functions in \eqref{eq:normalized-abel-jet} are holomorphic on a
neighborhood of $|s-1|\le r$. For every integer $k\ge0$,
\begin{equation}\label{eq:abel-jet-error}
 \left|\mathcal J_P^{(k)}(1)-\mathcal J_\lambda^{(k)}(1)\right|
 \le k!r^{-k}C_r\bigl(e^{\Delta(P,1-r)}-1\bigr).
\end{equation}
For each fixed $r$ and $K$, the right side tends to zero uniformly for
$0\le k\le K$ as $P\to\infty$.
\end{theorem}
\begin{proof}
Euler summation, first for $\R s>1$ and then by continuation to
$\R s>0$, gives
\[
 \zeta(s)=\frac{s}{s-1}-sI(s),\qquad
 I(s)=\int_1^\infty\{x\}x^{-s-1}\dd x.
\]
The integral is locally uniformly convergent in this half-plane and
$|I(s)|\le1/\R s$. On $|s-1|\le r<1/2$, therefore,
\begin{align*}
 |(s-1)\zeta(s)|
 &=|s|\,|1-(s-1)I(s)|\\
 &\ge(1-r)\left(1-\frac r{1-r}\right)=1-2r>0.
\end{align*}
The removable value at one is one. The numerator $\zeta(2s)$ is
holomorphic here and has absolute value at most $\zeta(2-2r)$ by
its absolutely convergent series. Hence
\[
 \sup_{|s-1|\le r}|\mathcal J_\lambda(s)|\le C_r.
\]
Since $1-r>\beta$, Lemma~\ref{lem:jet-euler-bound} bounds
$|G_P-1|$ on the same disk. Apply the Cauchy integral formula to
$\mathcal J_P-\mathcal J_\lambda$ on its boundary to obtain
\eqref{eq:abel-jet-error}. The exponent $\beta-1+r$ in
$\Delta(P,1-r)$ is negative, which proves convergence.
\end{proof}

\begin{corollary}[A completely explicit cutoff]\label{cor:jet-cutoff}
Fix $K\ge0$ and $\eta>0$. Set $\sigma=1-r$,
$d=\sigma-\beta>0$ and
$M=C_r\max_{0\le k\le K}k!r^{-k}$. Any integer $P$ with
\begin{equation}\label{eq:jet-cutoff}
 P\ge\max\left\{2,
 \left[\frac1{d(1-2^{-2\sigma})\log(1+\eta/M)}\right]^{1/d}
       \right\}
\end{equation}
satisfies $|\mathcal J_P^{(k)}(1)-\mathcal J_\lambda^{(k)}(1)|\le\eta$
for all $0\le k\le K$.
\end{corollary}
\begin{proof}
For $P\ge2$, \eqref{eq:jet-delta} is at most
$P^{-d}/[d(1-2^{-2\sigma})]$. The stated cutoff makes this at most
$\log(1+\eta/M)$. Substitute in \eqref{eq:abel-jet-error}.
\end{proof}

\begin{theorem}[Failure of the critical local energy for the same functions]
\label{thm:jets-do-not-bound-energy}
Put $B_P(X)=\sum_{n\ge1}f_P(n)e^{-n/X}$ and
\[
 \mathcal H_{B_P}(X)=\int_{X/2}^{2X}|B_P(Y)|^2\frac{\dd Y}{Y}.
\]
For every fixed $P$, the function $f_P$ has all the logarithmic
summatory and additive-twist savings proved in the preceding
construction. Nevertheless, for every $0<\varepsilon<2\beta-1$,
\begin{equation}\label{eq:jet-energy-failure}
 \mathcal H_{B_P}(X)\ne O(X^{1+\varepsilon}).
\end{equation}
This holds for the very same function whose Abel jets satisfy
\eqref{eq:abel-jet-error}, not for a different function at each scale.
\end{theorem}
\begin{proof}
The factorization \eqref{round26:eq:factorization} has a simple pole
at the real point $\beta$: $\zeta(s+1-\beta)$ supplies that pole,
the remaining factor $R_f$ is nonzero, and
$\zeta(2\beta)/\zeta(\beta)\ne0$.
For the denominator, the alternating zeta series shows
$\zeta(\beta)<0$ for $0<\beta<1$.

Suppose instead that the local energy is $O(X^{1+\varepsilon})$.
For $\sigma>(1+\varepsilon)/2$, Cauchy--Schwarz on each dyadic interval
gives
\begin{align*}
 \int_{2^j}^{2^{j+1}}|B_P(Y)|Y^{-\sigma}\frac{\dd Y}{Y}
 &\le(\log2)^{1/2}2^{-j\sigma}
            \mathcal H_{B_P}(2^j)^{1/2}\\
 &\ll2^{j((1+\varepsilon)/2-\sigma)}.
\end{align*}
The sum over $j\ge0$ converges locally uniformly to the right of
$(1+\varepsilon)/2$. The integral from zero to one is entire,
by exponential decay and $|f_P(n)|=1$. Thus
\[
 \int_0^\infty B_P(Y)Y^{-s}\frac{\dd Y}{Y}
\]
is holomorphic on $\R s>(1+\varepsilon)/2$. Initially for $\R s>1$
it equals $\Gamma(s)F_{f_P}(s)$ by absolute convergence. Continuation
would then remove its genuine pole at $\beta$, since
$(1+\varepsilon)/2<\beta$ and $\Gamma(\beta)\ne0$. This is a
contradiction. The logarithmic estimates follow from the earlier
nonnegative convolution $f_P=\lambda*h$ and require no RH input.
\end{proof}

\begin{figure}[tb]
\centering
\begin{tikzpicture}[>=Latex,x=8cm,y=1cm,font=\small]
 \fill[blue!6] (1,0) ellipse[x radius=.125,y radius=1];
 \draw[->] (.68,0)--(1.2,0) node[right] {$\R s$};
 \draw[red!70!black,thick] (.75,-.65)--(.75,.65);
 \node[below,red!70!black] at (.75,-.65) {pole $\beta$};
 \draw[blue!65!black,thick] (1,0) ellipse[x radius=.125,y radius=1];
 \fill[blue!65!black] (1,0) circle[x radius=.004,y radius=.032];
 \node[below] at (1,0) {$1$};
 \node[above,align=center] at (1,1.15)
  {local Abel data converge uniformly\\and all finite jets can be matched};
 \node[below,align=center] at (.94,-1.35)
  {$r<1-\beta$: the pole persists outside the disk\\
   and the critical energy estimate fails};
\end{tikzpicture}
\caption{The approximation and the pole coexist. The illustration uses
$\beta=3/4$ and $r=1/8$; the disk is drawn schematically.}
\label{fig:abel-jet-disk}
\end{figure}

For a numerical tolerance with an exact certificate, take
$\beta=3/4$, $r=1/8$, $K=3$ and $P=2^{320}$. Since
$\zeta(7/4)\le1+\int_1^\infty t^{-7/4}\dd t=7/3$,
we may replace $C_r$ by $112/27$. Thus $M\le114688/9$, and
\[
 \Delta(P,7/8)=\frac{8\,2^{-40}}{1-2^{-560}}<1,\qquad
 \frac{114688}{9}\frac{\Delta(P,7/8)}{1-\Delta(P,7/8)}<10^{-6}.
\]
The last inequality is rational arithmetic. The bound
$e^x-1\le x/(1-x)$ for $0\le x<1$ follows by comparing the
exponential series with the geometric series. The first four Abel
derivatives are therefore within $10^{-6}$ of those of Liouville,
and every value $f_P(n)$ up to $2^{320}$ agrees exactly with
$\lambda(n)$. The fixed function still violates every local-energy
bound $O(X^{1+\varepsilon})$ with $0<\varepsilon<1/2$.
The check \texttt{anc/abel-jet-bounds.py} verifies the displayed
rational constants; the infinite construction and its pole are
proved analytically above.

Exact agreement is different. The earlier coefficient-rigidity result
shows that equality of the normalized leading Abel coefficient in the
completely multiplicative sign class forces $f=\lambda$. The present
theorem does not contradict it: here $G_P(1)>1$ at every finite cutoff,
although $G_P(1)\to1$. Approximate equalities, at any fixed precision
and finite derivative order, cannot replace the exact arithmetic
condition. Nor does the exact condition itself prove an upper bound
for the actual Liouville energy.

\subsection{The unresolved upper estimate and the Mellin step}

For the actual Liouville function, the bound still needed is
\begin{equation}\label{eq:remaining-critical-local-bound}
 \mathcal H_B(X)\ll_\varepsilon X^{1+\varepsilon}
 \qquad\text{for every }\varepsilon>0\text{ and all }X\ge1.
\end{equation}
The dyadic argument in the last proof shows directly that this bound
would make the actual Mellin integral holomorphic on $\R s>1/2$.
It would coincide there with $\Gamma(s)\zeta(2s)/\zeta(s)$ and
exclude every right-hand zero. Conversely, Proposition~\ref{prop:smoothed}
under RH bounds $B(Y)$ by $Y^{1/2+\varepsilon/2}$ and gives
\eqref{eq:remaining-critical-local-bound} by integration. Hence this is
an equivalence, not a proved upper estimate.

In particular, using Plancherel on a line $\R s>1/2$ for the actual
$B$ requires the corresponding Mellin convergence or $L^2$ control
first. Meromorphic continuation of the quotient by itself supplies
neither. Moving a contour past a hypothetical zero also introduces
its residue; it does not discard the corresponding large mode.
The signed off-diagonal in \eqref{eq:local-off-diagonal} remains
unbounded at the required power. No new zero-free half-plane or
proof of RH is claimed in this revision.
